\clearpage
\section{Annotated term contexts and concrete closing observations}
\label{sec:copying-observations}
This section constructs the remaining syntax and observation interfaces of
Appendix~B.1, item~3, without claiming its congruence theorem. The Figure~16
translation of both proposition programs remains unchanged. Two different
composition operations must be distinguished: composing graph-building
constructors is proved here, whereas correspondence with addressed graph
plugging is still required. Execution remains an explicit parameter until
the copying UAM of item~2 is constructed.

\subsection{Exact single-hole closing contexts}
\termdef{copying-sole-hole}


A sole-hole certificate identifies the selected occurrence as the only hole
at every depth. At a shallow occurrence it requires a hole label there and
global hole-freedom below every other edge. Inside a box it requires the
corresponding certificate for the body and hole-freedom elsewhere outside.
Induction on the occurrence proves that filling it with a hole-free graph
leaves no holes. In the shallow case, every remaining edge is either a
retained non-hole or an inserted hole-free edge. In the deep case, only
the selected box body changes.

If the local input list is nonempty, the selected hole has a source: its
source-list length is its input-interface length. The same argument descends
through boxes. Thus the sourced-hole premise of the previous section follows
from the interface and sole-hole certificate; it is not a new assumption on
translated term contexts.

\lean{Copying/Contexts}{10}{19}{Contexts.SoleHole}


\lean{Copying/Contexts}{20}{41}{Contexts.fill_holeFree}


\lean{Copying/Contexts}{42}{63}{Contexts.sourced}


A closing context records an arbitrary outer depth, a graph of interface
$\star\Rightarrow\epsilon$, an addressed term hole of the requested depth
and environment interface, its sole-hole certificate, and binding-freedom.
The surrounding context can contain arbitrary example-language graph
structure permitted by these conditions; it is not restricted to syntax.
Filling is the existing addressed graph replacement. Hole-freedom follows
from the sole-hole lemma, focus-freedom from global replacement closure,
and binding-freedom from the sourced-hole theorem. The latter's source
condition is derived from the hole's single term input.

\lean{Copying/Contexts}{64}{71}{Contexts.ClosingContext}


\lean{Copying/Contexts}{72}{75}{Contexts.fill}


\lean{Copying/Contexts}{76}{79}{Contexts.fill_holeFree}


\lean{Copying/Contexts}{80}{83}{Contexts.fill_focusFree}


\lean{Copying/Contexts}{84}{88}{Contexts.fill_bindingFree}


\subsection{The annotated grammar and literal graph interpretation}
\termdef{copying-term-context}


A term context has a fixed hole scope and a separately indexed outer
scope. Its constructors are the hole, abstraction, left and right application,
unary operation, and left and right binary operation, exactly the grammar
of Section~11. An abstraction removes one variable from the outer scope;
the hole annotation retains any variables bound by enclosing abstractions.
Filling is structural recursion using the original term constructors.
Context composition substitutes an annotated context for the hole, with
matching scopes enforced by the type indices. Induction on the outer context
proves the filling law, the right identity law, and associativity. The left
identity is definitional. No variable names or program terms are changed.

\lean{Copying/TermContexts}{7}{17}{TermContexts.TermContext}


\lean{Copying/TermContexts}{18}{28}{TermContexts.fill}


\lean{Copying/TermContexts}{29}{39}{TermContexts.compose}


\lean{Copying/TermContexts}{40}{44}{TermContexts.fill_compose}


\lean{Copying/TermContexts}{45}{47}{TermContexts.compose_hole}


\lean{Copying/TermContexts}{48}{58}{TermContexts.compose_assoc}


A scoped graph packages its finite depth together with a graph of the
corresponding term and environment interface. This is a dependent pair, not a
quotient or a new translation. The packaged abstraction, unary, application,
and binary constructors call the existing Figure~16 constructors. Binary
branches use the original maximum-depth padding. Thus all distributor edges,
port ordering and boxes remain present.

Interpret a term context by replacing its hole with a supplied scoped graph
and applying these graph constructors recursively. Induction proves that
supplying the original translation of a term returns exactly the original
translation of the filled term, including its depth. Each inductive case
rewrites the recursive result and closes by definitional equality. A second
induction proves that interpretation respects context composition.

The actual holed translation is obtained by supplying the prescribed
single hole edge, at any chosen allowance and hole name. Its constructor
composition law follows from the preceding interpretation theorem. This
law does not yet identify interpretation with the operation that replaces
the selected edge in the holed graph: that addressed-plugging correspondence,
including padding and recursive isomorphism, remains to be proved.

\lean{Copying/TermContexts}{59}{60}{TermContexts.ScopedGraph}


\lean{Copying/TermContexts}{61}{62}{TermContexts.ofTerm}


\lean{Copying/TermContexts}{63}{65}{TermContexts.abstraction}


\lean{Copying/TermContexts}{66}{68}{TermContexts.unary}


\lean{Copying/TermContexts}{69}{73}{TermContexts.application}


\lean{Copying/TermContexts}{74}{80}{TermContexts.binary}


\lean{Copying/TermContexts}{81}{92}{TermContexts.apply}


\lean{Copying/TermContexts}{93}{105}{TermContexts.apply_ofTerm}


\lean{Copying/TermContexts}{106}{112}{TermContexts.apply_compose}


\lean{Copying/TermContexts}{113}{115}{TermContexts.hole}


\lean{Copying/TermContexts}{116}{118}{TermContexts.graph}


\lean{Copying/TermContexts}{119}{123}{TermContexts.graph_compose}


\subsection{Termination observations with an explicit execution interface}
\termdef{copying-closing-observation}


An open program contains an actual graph of a fixed term/environment
interface and proofs of global focus-freedom and hole-freedom. A closed
program contains a graph of interface $\star\Rightarrow\epsilon$ with
the same two proofs and its finite depth. Closing an open program in a
certified context constructs a closed program using the filling lemmas.
An execution interface supplies a machine and an initializer for these
closed programs. The interface is a parameter: no copying machine, machine
law, or translation-correctness hypothesis is asserted by its definition.

\lean{Copying/Observations}{11}{15}{Observations.OpenProgram}


\lean{Copying/Observations}{16}{21}{Observations.ClosedProgram}


\lean{Copying/Observations}{22}{27}{Observations.ClosingContext.close}


\lean{Copying/Observations}{28}{35}{Observations.Execution}


An observation interface fixes the fragment's depth and the two environment
sizes. Its context type is the concrete closing-context structure above,
whose outer depth is unrestricted. Observation initializes the filled closed
graph. Since validity is built into that structure, every inhabitant is
allowed. Successful termination means a finite machine run to a final state;
a stuck nonfinal state is not counted.

Quantitative refinement universally quantifies over these concrete contexts
and compares their initialized states by the existing machine refinement.
Equivalence is mutual refinement. Unfolding the observation relation proves
that this direct definition equals the generic observation interface used
by the robustness infrastructure. Reflexivity and transitivity hold pointwise
from the corresponding machine results; equivalence is consequently reflexive,
symmetric and transitive. No context closure is used to prove these order
laws. Closure is needed for congruence, a distinct theorem still outstanding.

For the universal cost relation, refinement is exactly implication of
successful termination in each context. One direction extracts the target
run promised by refinement; the other supplies that run and the trivial
cost witness. Applying both directions proves that equivalence is precisely
agreement on termination, without determinism or a runtime bound.

\lean{Copying/Observations}{36}{40}{Observations.Interface}


\lean{Copying/Observations}{41}{47}{Observations.observations}


\lean{Copying/Observations}{48}{52}{Observations.Terminates}


\lean{Copying/Observations}{53}{56}{Observations.Refines}


\lean{Copying/Observations}{57}{61}{Observations.Equiv}


\lean{Copying/Observations}{62}{71}{Observations.refines_iff_observations}


\lean{Copying/Observations}{72}{74}{Observations.refines_refl}


\lean{Copying/Observations}{75}{78}{Observations.refines_trans}


\lean{Copying/Observations}{79}{81}{Observations.equiv_refl}


\lean{Copying/Observations}{82}{84}{Observations.equiv_symm}


\lean{Copying/Observations}{85}{90}{Observations.equiv_trans}


\lean{Copying/Observations}{91}{102}{Observations.refines_any_iff}


\lean{Copying/Observations}{103}{114}{Observations.equiv_any_iff}


The original translation of a term is packaged at any adequate allowance
using the existing padding operation. Its global label theorems prove
focus-freedom and hole-freedom after padding. Term observational equivalence
then compares both translations at every common adequate allowance, against
all certified closing contexts, using the universal cost relation. The
execution parameter remains visible: the concrete copying UAM must instantiate
it later. Padding/initialization invariance and the full paper-correspondence
argument are not consequences of this definition alone.

\lean{Copying/Observations}{115}{122}{Observations.translatedProgram}


\lean{Copying/Observations}{123}{127}{Observations.ObservationalEquiv}


The closing-context type is nonempty for closed terms at every allowance:
the identity observation is the single $\star\Rightarrow\epsilon$ hole.
Its sole-hole proof uses the singleton edge carrier and its binding-freedom
is the isolated-hole theorem. This checks that the observation quantifier
is not vacuous in the closed case.

Three exact translation checks exercise the annotation and constructor laws.
An enclosing lambda binds the variable at index zero in the hole's annotated
scope. The left proposition program is reconstructed by filling the argument
of its original lambda application with $\mathtt{ref}\,1$. The right program
is reconstructed by filling an abstraction context with its original sequence
$f();1$. Each proof invokes the exact interpretation theorem. These are
structural equalities, not semantic equivalence claims.

\lean{Copying/ObservationExamples}{10}{24}{ObservationExamples.identityClosingContext}


\lean{Copying/ObservationExamples}{25}{32}{ObservationExamples.bound_hole_translation}


\lean{Copying/ObservationExamples}{33}{43}{ObservationExamples.privateCell_context_factor}


\lean{Copying/ObservationExamples}{44}{49}{ObservationExamples.constantResult_context_factor}


The full build and three axiom audits pass. Binding-freedom is now proved in
Section~\ref{sec:copying-context-binding}. The sole-hole certificate is proved in
Section~\ref{sec:copying-certified-contexts}. Item~3 remains open pending addressed
plugging correspondence and term-context congruence with matching observation
and padding quantification. The execution interface
must also receive its copying-UAM implementation from item~2. No conditional
wrapper is presented as Proposition~12.2 or Lemma~11.3.
